\begin{abstract}
Η αυτόματη τμηματοποίηση του πλακούντα σε μαγνητική τομογραφία (\en{MRI}) αποτελεί θεμελιώδες βήμα
για την ποσοτική ανάλυση της κύησης και τη συστηματική αξιοποίηση πληροφοριών
ιατρικής απεικόνισης. Η διαδικασία είναι ιδιαίτερα απαιτητική, λόγω έντονης
μεταβλητότητας στη μορφολογία και στην εμφάνιση του οργάνου, ασαφών ορίων, και
της περιορισμένης διαθεσιμότητας επαρκώς επισημασμένων δεδομένων.

Η παρούσα διπλωματική εργασία πραγματοποιεί συστηματική συγκριτική αξιολόγηση
σύγχρονων αρχιτεκτονικών βαθιάς μάθησης για \en{3D} τμηματοποίηση πλακούντα.
Αναπτύχθηκε ενιαίο, αναπαραγώγιμο πειραματικό πλαίσιο σε \en{MONAI/PyTorch},
το οποίο ενοποιεί προεπεξεργασία, δειγματοληψία, επαύξηση, εκπαίδευση και
αξιολόγηση, ώστε η σύγκριση να είναι δίκαιη και τεκμηριωμένη.

Η μελέτη καλύπτει αρχιτεκτονικές τύπου \en{U-Net}, \en{Transformer}-βασισμένες
προσεγγίσεις και \en{state-space}-βασισμένα μοντέλα, με αποτίμηση μέσω
\en{Dice/IoU} και ποιοτικής επιθεώρησης των παραγόμενων μασκών. Τα ευρήματα
αναδεικνύουν ότι η αυστηρή τυποποίηση του \en{pipeline} αποτελεί κρίσιμο
παράγοντα για αξιόπιστα συμπεράσματα και ότι διαφορετικές οικογένειες
αρχιτεκτονικών παρουσιάζουν διακριτά πλεονεκτήματα ως προς τη σταθερότητα, την
ποιότητα ορίων και τη συνολική συμπεριφορά τμηματοποίησης.

Συνολικά, η εργασία προσφέρει ένα σαφές πλαίσιο αναφοράς για την επιλογή
μοντέλων τμηματοποίησης πλακούντα σε \en{MRI}. Παράλληλα, υπογραμμίζει την ανάγκη
για περαιτέρω διερεύνηση της γενίκευσης σε ανεξάρτητα σύνολα δεδομένων, για
εξωτερική επικύρωση, καθώς και για ενσωμάτωση τεχνικών εκτίμησης αβεβαιότητας.

   \begin{keywords}
   \en{MRI}, Πλακούντας, Τμηματοποίηση ιατρικών εικόνων, Βαθιά Μάθηση, \en{MONAI},
   Συγκριτική αξιολόγηση, \en{U-Net}, \en{Transformers}, \en{State Space Models},
   Αναπαραγωγιμότητα
   \end{keywords}
\end{abstract}
